% Pursuing Stacks (cotes 134-2 to 134-8): archive cote 134-7, site page 102 (PDF page 103), typescript page 544.
% typescript. Draft compared: ps6.tex:5133-5196.
% AI-transcribed by Claude Opus 5.5 (claude-opus-5-5, adaptive thinking, effort high), Message Batches API, batch ps4 (msgbatch_012Aoc1NuskATPsyBeKSYK7J), 2026-11-12; not yet checked by eye.
\scanpage{134-7/102}
\mspage{54\over{\illegible{1}}{4}}

130.\qquad In a more down-to-earth line,  and \xout{\unsure{still}} reverting to the
``unipotent'' approach still (in either W-or V-version for ``unipotent bund-%
les''),  I would still like to point out one rather mild extension of the
set-up as contemplated initially.  The sugge stion here is to admit as
compnents for our semisimplicial models not merely sheaves of sets which
are ``unipotent bundles'',  but equally \xout{finite} direct sums of such.  \xout{The x}
\inserted[typed]{After making such extension,} \xout{\illegible{2} \unsure{restriction is needed if we want a}} \inserted{the} linearization functor L
((11) p. 539) \xout{\unsure{is}} \inserted{still} makes a sense\hand{,} \deleted{presumably,}  provided $Uab(k)$ is stable
under infinite sums (otherwise, we'd have to restrict to a finite number
of connected components for our ``bundles'').  Presumably,  once we get
into this,  we will have to admit ``twisted'' finite direct sums as well,
in order to have \inserted{basic} notions compatible with descent  - never mind such
technicalities at the present stage of reflections !  Thus,  applying
componentwise (i.e. to each component $X_n$) the $\hand{\pi}_o$-functor (``connected
components''),  a semisimplicial ``bundle'' $X_{*}$ in the wider sense give\hand{s} rise
to an associated  usual semisimplicial set,  let's call it $\hand{\pi}_{o/k}(X_{*})$,
together with a map
\[ (18) \qquad X_{*} \hand{\longrightarrow} \hand{\pi}_{o/k}(X_{*}) \quad , \]
where the second hand side is interpreted as a semisimplicial constant
object over k.  (For simplicity,  we have assumed here Spec(k) to be
0-connected,  and that the direct sums involved in the $X_n$'s  are not
twisted ...)Intuitively,  we may interprete (18) as defining $X_{*}$ as a
(``strict'', namely componentwise connected) \uline{schematic} homotopy type,
lying ``over'' the \uline{discrete} (or ``constant'') homotopy type $\hand{\pi}_{o/k}(X_{*}) = E_{*}$ .
The \over{\illegible{1}}{l}atter introduces the homotopy invariants of it's own,  which strict
-ly speaking should'nt be vi\inserted{e}wed as  being of a ``schematic'' nature.  Thus,
when our exclusive emphasis is on studying the  ``strict'' schematic
homotopy types,  we'll restrict our models $X_{*}$ by demanding that the
associated discrete homotopy type $\hand{\pi}_{o/k}(X_{*})$  should be \uline{aspheric},  i.e.
isomorphic in Hot to a one-point space.  Under this restriction,  pre-%
sumably,  working with those slightly more general models should give
(up to equivalence) the same derived category $Hot_o(k)$,  as when working
with (connected) unipotent bundles.  \xout{The usefulness of} Allowing connec-%
ted components may prove useful for giving a little more ``elbow freedom''
in working with models,  as it allows for instance anodine operations
such as taking direct sums.  Thus,  all constant semisimplicial objects
will be allowed,  or at any rate those which correspond to aspheric semi
simplicial sets  -  which includes notably the objects of \hand{$\Delta^{\wedge}$} represen-%

\margin{102}
